\section{Warmup: Simply increasing the module rank}
\label{sec:warmup}
The main mechanism to increase the security margin in \textsc{Kyber} is to raise
the module rank $k$~\cite{kyberround1}. Fixing the ML-KEM base ring, noise distribution, and
modulus, and varying only $k$, we measure the decryption-failure
probability $\delta$ using the \textsc{Kyber} team's scripts~\cite{kyber-scripts}.
Throughout we use the centered binomial noise $\mathrm{CBD}_2$
(the noise of ML-KEM-768 and -1024; recall that ML-KEM-512 uses $\eta_1=3$).
The ranks $k=3 \text{ and } 4$ recover the uncompressed analogue of standardized sets ML-KEM-768, and -1024,
while $k \geq 5$ is an extrapolation.

We require $\delta\le 2^{-128}$ throughout the paper. This is the same bar as for
\textsc{Kyber}, whose parameter sets all achieve a failure probability below $2^{-128}$
with some margin~\cite{kyberround3}. Since $\delta$ is information-theoretic in the
classical random-oracle model, it does not need to decrease with the security
parameter. Here $\delta$ must be the failure
probability in the adversarial sense used by Fujisaki--Okamoto proofs, in which the
adversary knows the secret key and chooses the message. The scripts instead compute the
statistical failure probability of Section~\ref{sec:preliminaries}, that is, the probability
over the randomness of key generation and encryption that some coefficient of $\tilde{n}$
exceeds $q/4$. Without ciphertext compression, as in all experiments in this paper,
$\tilde{n}$ does not depend on the message, so the adversarial failure probability is
bounded by the statistical one. We machine-check this in EasyCrypt
(Section~\ref{sec:verification}) and report the resulting bounds alongside the script
estimates in Table~\ref{tab:cbd-baseline}.

Against this target, the bound holds at $k=7$ ($2^{-133.7}$ by the scripts, $2^{-133}$
machine-checked) but fails at $k=8$ ($2^{-116.8}$ and $2^{-116}$, respectively), which demonstrates that raising
$k$ alone cannot buy higher security without sacrificing correctness.

With \textsc{Kyber}'s ciphertext compression, the rounding of $\mathbf{u}$ and $v$ would add
further noise, so the correctness limit on $k$ would be reached even earlier.

\begin{table}[h]
\centering
\caption{Decryption-failure probability and public-key and ciphertext
sizes as a function of module rank $k$, at fixed ML-KEM noise and modulus,
without ciphertext compression. Ranks $k=3,4$ correspond to
uncompressed ML-KEM-768, and -1024. \emph{Script}: estimate from the \textsc{Kyber} team's
scripts. \emph{Machine-checked}: upper bound on the adversarial failure probability proved in
EasyCrypt (Section~\ref{sec:verification}), with the final tail probability computed by OCaml
code that rounds upward.}
\label{tab:cbd-baseline}
\begin{tabular}{ccccc}
\toprule
 & \multicolumn{2}{c}{Failure probability} & & \\
\cmidrule(lr){2-3}
$k$ & Script & Machine-checked & pk (bytes) & ct (bytes) \\
\midrule
3 & $2^{-299.6}$ & $2^{-300}$ & 1184  & 1536 \\
4 & $2^{-231.1}$ & $2^{-230}$ & 1568  & 1920 \\
5 & $2^{-186.5}$ & $2^{-186}$ & 1952  & 2304 \\
6 & $2^{-155.9}$ & $2^{-155}$ & 2336  & 2688 \\
7 & $2^{-133.7}$ & $2^{-133}$ & 2720  & 3072 \\
\rowcolor{failbg}
8 & $2^{-116.8}$ & $2^{-116}$ & 3104  & 3456 \\
\bottomrule
\end{tabular}
\end{table}
